% Generated draft: requires mathematical and source review.
\documentclass[11pt]{article}
\usepackage[margin=25mm]{geometry}
\usepackage{fontspec}
\setmainfont{Latin Modern Roman}
\usepackage{amsmath,amssymb,mathtools,unicode-math}
\setmathfont{Latin Modern Math}
\usepackage{enumitem,needspace,xurl,hyperref}
\setlength{\parindent}{0pt}
\setlength{\parskip}{5pt}
\setlength{\emergencystretch}{3em}
\newcommand{\N}{\mathbb{N}}
\newcommand{\Z}{\mathbb{Z}}
\newcommand{\Q}{\mathbb{Q}}
\newcommand{\R}{\mathbb{R}}
\newcommand{\C}{\mathbb{C}}
\newcommand{\F}{\mathbb{F}}
\newcommand{\A}{\mathbb{A}}
\newcommand{\PP}{\mathbb{P}}
\newcommand{\E}{\mathbb{E}}
\newcommand{\eps}{\varepsilon}
\newcommand{\dd}{\,\mathrm{d}}
\begin{document}
\section*{UnsolvedMath ID: 2307064}
\textbf{Source problem code:} catalogue entry
\subsection{Definitions and mathematical statement}
Let $\mathbb D=\{z\in\C:|z|<1\}$, let $\Gamma$ be a Fuchsian group of holomorphic automorphisms of $\mathbb D$, and let $i(z)=z$ be its identity. The source does not assume torsion-freeness. For a disk automorphism $\gamma$, $|\gamma'(0)|=1-|\gamma(0)|^2$.

For a literal nonnegative interpretation, define the sum in the extended sense and the product as the limit over finite subsets:
\[
 A(\Gamma)=\sum_{\gamma\in\Gamma}|\gamma'(0)|\in[0,\infty],\qquad
 B(\Gamma)=\prod_{\substack{\gamma\in\Gamma\\\gamma\ne i}}|\gamma(0)|^2\in[0,1].
\]
A zero factor makes $B=0$, and a divergent nonnegative sum is assigned $+\infty$. If only positive convergent products or finite sums are intended, extra hypotheses are missing.

\textbf{Recorded question (N. Suita, communicated by Ch. Pommerenke).} Is
\[
 \sum_{\gamma\in\Gamma}|\gamma'(0)|\ \geq\
 \prod_{\substack{\gamma\in\Gamma\\\gamma\ne i}}|\gamma(0)|^2
\]
true? The statement says this is an equivalent formulation of a problem connecting the Bergman kernel and capacity function of a Riemann surface.

There is a substantive transcription/formulation concern: the left side includes the identity contribution $|i'(0)|=1$, while each factor on the right lies in $[0,1)$. Thus under the displayed nonnegative conventions the bound is automatic and cannot by itself identify the substantive kernel--capacity comparison alluded to. The exact absolute-value placement, normalization, and intended group-series expression must be checked; no alternative signed, complex, or renormalized formula is invented.

The registry classification is SOLVED-IN-LITERATURE, but its research summary itself labels the precise situation partial/cautionary: it says the underlying Suita conjecture has been proved and the specific normalized Fuchsian-group form still needs checking. The summary's description refers to a boundary point, whereas the displayed expression is evaluated at the interior point $0$ of $\mathbb D$ and no kernel or capacity definitions are provided. The intended correspondence and point conventions require review. The solved label is therefore not treated as authentication of this exact formula.
\subsection{Short English statement}
The recorded Fuchsian-group inequality is attributed to Suita and linked to a Bergman-kernel/capacity problem. Its literal absolute-value form is automatic from the identity term and unit-disc factors, so the intended substantive formula needs checking. The summary expressly cautions that the exact group formulation is unverified despite the solved Suita-conjecture label.
\subsection{Sources}
\begin{itemize}
\item Local UnsolvedMath registry, permanent ID 2307064, source code catalogue entry, ``Research Problems in Function Theory --- Problem 7.64''; attribution: N. Suita; communicated by Ch. Pommerenke. Supplied status solved, classification SOLVED-IN-LITERATURE, updated 2026-08-06; the research summary explicitly qualifies the exact formulation as needing review.
\item Hayman, \emph{Research Problems in Function Theory} (2018), Problem 7.64, \url{https://arxiv.org/abs/1809.07200}. The record reports source-TeX extraction and access on 2026-07-29; the 2018 update was retained for status triage. Current status and rights require review; the record describes publicly available arXiv source TeX.
\item The supplied summary reports proof of the underlying Suita conjecture but gives no proof reference or precise identification of this Fuchsian-group formula with that theorem.
\end{itemize}
\end{document}
